\documentclass{article}

\usepackage[margin=1in]{geometry}
\usepackage{hyperref}
\usepackage{listings}
\usepackage{graphicx}
\usepackage{booktabs}
\usepackage{enumitem}
\usepackage{xcolor}

\hypersetup{
  colorlinks=true,
  linkcolor=blue,
  urlcolor=blue,
  pdftitle={Level 2 -- RAG-Enabled Chatbot (Elementary)}
}

\lstset{
  basicstyle=\ttfamily\small,
  breaklines=true,
  frame=single,
  backgroundcolor=\color{gray!10}
}

\title{Level 2 -- RAG-Enabled Chatbot (Elementary)}
\author{Swaroop Formulation Industries}
\date{2026-03-25}

\begin{document}
\maketitle

\section{Goal}

Extend the chatbot so it \textbf{retrieves answers from ingested documents} using \textbf{vector search}, not only parametric model knowledge.

\section{New Technology}

\begin{tabular}{@{}lp{8.5cm}@{}}
\toprule
\textbf{Component} & \textbf{Role} \\
\midrule
\textbf{ChromaDB} & Local, free vector database \\
\textbf{HuggingFace} & \texttt{sentence-transformers/all-MiniLM-L6-v2} embeddings \\
\textbf{LangChain} & PDF loading, chunking, RAG chain composition \\
\textbf{FastAPI} & HTTP API for chat and ingestion triggers \\
\bottomrule
\end{tabular}

\section{Features}

\begin{itemize}[leftmargin=*]
  \item \textbf{Ingest} the project report PDF and \textbf{ISO certificates} into the vector store.
  \item Pipeline: \textbf{Query $\rightarrow$ Embed $\rightarrow$ Retrieve top-k $\rightarrow$ Augment prompt $\rightarrow$ Generate}.
  \item \textbf{\texttt{/api/chat}} endpoint with \textbf{streaming} responses.
  \item \textbf{Source citations} in the UI (chunk metadata / filenames).
  \item \textbf{About} page describing data sources and limitations.
\end{itemize}

\section{Data Ingestion Script (Python)}

Example pattern using LangChain loaders, splitter, Chroma, and HuggingFace embeddings:

\begin{lstlisting}[language=Python]
from langchain_community.document_loaders import PyPDFLoader
from langchain_text_splitters import RecursiveCharacterTextSplitter
from langchain_community.vectorstores import Chroma
from langchain_community.embeddings import HuggingFaceEmbeddings

loader = PyPDFLoader("path/to/document.pdf")
docs = loader.load()

splitter = RecursiveCharacterTextSplitter(
    chunk_size=1000,
    chunk_overlap=200,
)
chunks = splitter.split_documents(docs)

embeddings = HuggingFaceEmbeddings(
    model_name="sentence-transformers/all-MiniLM-L6-v2"
)

vectorstore = Chroma.from_documents(
    documents=chunks,
    embedding=embeddings,
    persist_directory="./chroma_db",
)
vectorstore.persist()
\end{lstlisting}

Adjust paths, collection names, and chunk sizes for your corpus.

\section{Implementation Steps}

\begin{enumerate}[leftmargin=*]
  \item \textbf{Set up} a FastAPI backend project structure and dependencies.
  \item \textbf{Create} an ingestion script (or CLI) for PDFs and certificate files.
  \item \textbf{Build} the RAG chain with LangChain (retriever + LLM binding).
  \item \textbf{Add} a streaming \textbf{\texttt{/api/chat}} endpoint (SSE or chunked responses).
  \item \textbf{Update} the frontend to call the backend instead of only Groq from the client.
  \item \textbf{Add} an \textbf{About} page documenting ingestion, models, and privacy.
  \item \textbf{Deploy} the backend to \textbf{Render} (or comparable) with persistent disk or rebuild-on-deploy strategy for Chroma.
\end{enumerate}

\section{Deployment Options}

\begin{tabular}{@{}lp{7.5cm}@{}}
\toprule
\textbf{Platform} & \textbf{Notes} \\
\midrule
\textbf{Render.com} & Free tier suitable for hobby demos; verify sleep/cold start behavior. \\
\textbf{HuggingFace Spaces} & Good for demos with Gradio/Streamlit-style wrappers. \\
\textbf{Railway.app} & Paid from about \textbf{\$5/month} for more consistent uptime. \\
\bottomrule
\end{tabular}

\section{Environment Variables}

\begin{tabular}{@{}lp{7cm}@{}}
\toprule
\textbf{Variable} & \textbf{Purpose} \\
\midrule
\texttt{CHROMA\_PERSIST\_DIR} & Directory for persisted Chroma data \\
\texttt{EMBEDDING\_MODEL} & HuggingFace model id (e.g.\ \texttt{sentence-transformers/all-MiniLM-L6-v2}) \\
\texttt{GROQ\_API\_KEY} & LLM inference (or substitute provider) \\
\bottomrule
\end{tabular}

\section{Version Control}

\begin{itemize}[leftmargin=*]
  \item Use a feature branch such as \textbf{\texttt{feature/rag-v1}} for RAG work.
  \item Tag milestone: \textbf{\texttt{v0.2.0-l2-rag-enabled}}.
\end{itemize}

\section{Python Requirements}

Install (versions pinned in your own \texttt{requirements.txt}):

\begin{itemize}[leftmargin=*]
  \item \texttt{langchain}
  \item \texttt{chromadb}
  \item \texttt{sentence-transformers}
  \item \texttt{fastapi}
  \item \texttt{uvicorn}
  \item \texttt{pypdf}
\end{itemize}

Also include LangChain community / text-splitter packages as needed for your import paths (\texttt{langchain-community}, \texttt{langchain-text-splitters}, etc.).

\end{document}
